\section{Genealogical moments, Jacobi operators, and spectral residues}
\label{sec:jacobi-generated}

The calendar-induced measure provides a link between cell geometry and
a positive spectral realization. This link uses the entire marked
partition: each of its twelve intervals preserves the same chronological
mass, whereas its length comes from the corresponding register
coordinate. The density exactly compensates for that length.
Consequently, the spatial distribution of the measure retains the
anisotropy of the register even though the distribution among windows
is uniform.

Let \(K_j>0\), \(Q=\sum_{j=1}^{12}K_j\), \(d_j=K_j/Q\), and
\(t_j=\sum_{i\leq j}d_i\), with \(t_0=0\). The probability measure
used in this section is
\[
 \dd\mu_K(t)=\sum_{j=1}^{12}
 \frac{\mathbf1_{[t_{j-1},t_j]}(t)}{12d_j}\,\dd t.
\]
Shared endpoints have measure zero. Its moments are rational for
every integer register:
\begin{equation}
 m_r=\int_0^1t^r\,\dd\mu_K(t)
 =\frac1{12(r+1)}\sum_{j=1}^{12}
   \frac{t_j^{r+1}-t_{j-1}^{r+1}}{d_j},\qquad r\geq0.
 \label{eq:jacobi-moments}
\end{equation}
The equality \(m_0=1\) expresses conservation of total mass. For
each \(n\geq1\), the Hankel matrix \(H_n=(m_{a+b})_{0\leq a,b<n}\)
is positive definite: its quadratic form is the integral of the square
of a nonzero polynomial over an interval where the density is positive.
This observation allows passage from moments to an operator without
choosing a target spectrum.

\begin{theorem}[Tridiagonal recurrence determined by the partition]
\label{thm:jacobi-recurrence}
There is a unique sequence of monic polynomials \(p_n\), of degree
\(n\), orthogonal for \(\mu_K\). If
\[
 h_n=\int p_n^2\,\dd\mu_K>0,\qquad
 a_n=\frac{\int t p_n^2\,\dd\mu_K}{h_n},\qquad
 \beta_n=\frac{h_n}{h_{n-1}}>0\quad(n\geq1),
\]
then \(p_0=1\), \(p_1=t-a_0\), and
\begin{equation}
 p_{n+1}(t)=(t-a_n)p_n(t)-\beta_np_{n-1}(t),\qquad n\geq1.
 \label{eq:jacobi-recurrence}
\end{equation}
All coefficients \(a_n,\beta_n\) are rational for the stated integer
register. In the orthonormal basis \(\psi_n=p_n/\sqrt{h_n}\),
the compression of multiplication by \(t\) to polynomials of degree
less than \(n\) has the symmetric matrix
\[
 J_n=\begin{pmatrix}
 a_0&\sqrt{\beta_1}&&0\\
 \sqrt{\beta_1}&a_1&\ddots&\\
 &\ddots&\ddots&\sqrt{\beta_{n-1}}\\
 0&&\sqrt{\beta_{n-1}}&a_{n-1}
 \end{pmatrix}.
\]
\end{theorem}
\begin{proof}
Positivity of \(H_n\) guarantees that the linear system imposing
orthogonality on a monic polynomial has a unique solution. Since its
coefficients are rational moments, the solution is rational. Expand
\(tp_n\) in the monic basis through degree \(n+1\). For
\(j\leq n-2\), symmetry of the inner product gives
\(\langle tp_n,p_j\rangle=\langle p_n,tp_j\rangle=0\).
Only the terms of degrees \(n+1,n,n-1\) remain. Their coefficients
are, respectively, \(1,a_n,h_n/h_{n-1}\), proving the recurrence.
Normalization by \(\sqrt{h_n}\) turns the two adjacent coefficients
into \(\sqrt{\beta_n}\).
\end{proof}

For the register expressed in Section~\ref{sec:register}, the first
diagonal entry and the square of the first coupling are
\[
 a_0=m_1=\frac{71195}{150312},\qquad
 \beta_1=m_2-m_1^2=\frac{2206226167}{22593697344}.
\]
The first equality calculates the center of mass of the chronological
partition; the second calculates its dispersion and determines the
first link of the tridiagonal matrix. Both numbers are obtained by
integrating the generated lengths, and their structural meaning
precedes any decimal evaluation.

\begin{theorem}[Simple spectrum, quadrature, and positive residues]
\label{thm:jacobi-residues}
The eigenvalues \(\lambda_1<\cdots<\lambda_n\) of \(J_n\) are
simple and belong to \((0,1)\). There are unique weights \(w_i>0\),
summing to one, such that
\begin{equation}
 \int f(t)\,\dd\mu_K(t)=\sum_{i=1}^{n}w_i f(\lambda_i)
 \qquad(\deg f\leq2n-1).
 \label{eq:jacobi-quadrature}
\end{equation}
The resolvent observed at the unit constant \(e_0\) satisfies
\begin{equation}
 g_n(z)=\langle e_0,(zI-J_n)^{-1}e_0\rangle
       =\sum_{i=1}^{n}\frac{w_i}{z-\lambda_i}
       =\cfrac1{z-a_0-\cfrac{\beta_1}{z-a_1-
          \cfrac{\beta_2}{\ddots-\cfrac{\beta_{n-1}}{z-a_{n-1}}}}}.
 \label{eq:jacobi-resolvent}
\end{equation}
In particular, its poles and residues are obtained from the register
through the operations of taking moments, orthogonalization, and
compression.
\end{theorem}
\begin{proof}
For any nonzero polynomial \(p\) of degree less than \(n\),
\[
 0<\int_0^1 t|p(t)|^2\,\dd\mu_K(t)
   <\int_0^1 |p(t)|^2\,\dd\mu_K(t).
\]

The variational characterization places the spectrum of \(J_n\) in
\((0,1)\). Each subdiagonal entry is positive. The eigenvector
equations recursively determine all its coordinates from the first;
if that coordinate vanished, the entire vector would vanish. Every
eigenspace therefore has dimension one. Symmetry of the matrix gives
orthogonal diagonalization and simplicity of the spectrum. For unit
eigenvectors \(v_i\), set \(w_i=|\langle e_0,v_i\rangle|^2>0\).
Completeness of the eigenbasis gives \(\sum_iw_i=1\) and the first
expression for the resolvent.

By expansion of the determinant, the tridiagonal recurrence shows that
\(\det(tI-J_n)=p_n(t)\). For \(\deg f\leq2n-1\), Euclidean
division writes \(f=qp_n+r\), with \(\deg q,\deg r<n\).
Orthogonality makes the integral of \(qp_n\) vanish, and its
evaluation at the \(\lambda_i\) also vanishes. For \(r\) of degree
less than \(n\), successive multiplications of the constant by \(t\)
remain within the polynomial space up to the required degree; hence
\(\int r\,\dd\mu_K=\langle e_0,r(J_n)e_0\rangle\).
The spectral resolution proves the quadrature formula. Its uniqueness
follows from invertibility of the Vandermonde matrix evaluated at the
distinct roots. Finally, the Schur complement of the first row and
column of \(zI-J_n\), iterated over the trailing submatrices,
produces the continued fraction.
\end{proof}

\begin{figure}[htbp]
\centering
\begin{tikzpicture}[>=Latex, node distance=8mm and 8mm,
 every node/.style={font=\small}, box/.style={draw,rounded corners=1pt,
 text width=37mm,minimum height=13mm,align=center}]
\node[box] (k) {Marked register\\\(K,Q\)};
\node[box,right=of k] (m) {Chronological measure\\\(\mu_K\)};
\node[box,right=of m] (h) {Positive moments\\\(H_n\)};
\node[box,below=of h] (j) {Tridiagonal operator\\\(J_n\)};
\node[box,left=of j] (g) {Resolvent\\\(g_n(z)\)};
\node[box,left=of g] (w) {Poles and residues\\\(\lambda_i,w_i>0\)};
\draw[->] (k)--(m);\draw[->] (m)--(h);\draw[->] (h)--(j);
\draw[->] (j)--(g);\draw[->] (g)--(w);
\end{tikzpicture}
\caption{Spectral realization of the HMT partition. Cell geometry
determines the measure; positivity of its moments produces the operator
and the positive residues of its resolvent.}
\label{fig:jacobi-generated}
\end{figure}

The resolvent formula provides a spectral reading of spatial memory.
The positive weights describe how much the constant coordinate observes
of each mode; the poles are the algebraic frequencies of the compressed
multiplication operator. Moments of increasing order recover successive
details of the measure, and the continued-fraction coefficients preserve
this information through a local three-term recurrence. The classical
appearance of a Jacobi matrix is recognized after its coefficients have
been constructed from the HMT cells. The causal chain is fixed by
\[
 K\longmapsto(d,t)\longmapsto\mu_K\longmapsto(m_r)_{r\geq0}
 \longmapsto\bigl((a_n)_{n\geq0},(\beta_n)_{n\geq1}\bigr)
 \longmapsto(J_n,g_n).
\]

\begin{proposition}[Dimension and refinement compatibility]
\label{prop:jacobi-refinement}
The matrices \(J_n\) are compatible principal compressions of the
same bounded multiplication by \(t\) on \(L^2(\mu_K)\). Their
spectral measures \(\sum_iw_i\delta_{\lambda_i}\) converge weakly
to \(\mu_K\). If each nonadic subinterval is represented by its
midpoint and its mass \(1/(12\cdot9^r)\), the resulting atomic
measure \(\mu_{K,r}\) approximates the moments of \(\mu_K\).
For a fixed degree \(n\), construction from these moments recovers
the coefficients \(a_j,\beta_j\) of the same matrix as refinement
tends to infinite depth.
\end{proposition}
\begin{proof}
Orthogonalization of the complete sequence uses the same formulas as
\(n\) increases, so each principal block remains fixed. Quadrature
exactly reproduces any polynomial once \(2n-1\) exceeds its degree.
Uniform approximation of continuous functions on \([0,1]\) by
polynomials, together with total mass one, proves weak convergence.
The midpoint sums converge on each of the twelve intervals with its
constant density, and thus approximate all moments. In fixed dimension,
the orthogonalization coefficients are rational functions of finitely
many moments, with denominators formed from positive Gram determinants.
Convergence of these moments and positivity of the limit ensure
convergence of the coefficients.
\end{proof}

\begin{falsifier}
A zero length invalidates the stated density and can destroy the rank
of the moments. Replacing the fixed measure by unnormalized counts on
different meshes changes the operator. The residues of \(g_n\) are
spectral weights of the resolvent, whereas the residues of a canonical
form are defined on geometric divisors: identifying these notions
requires specifying the form and the corresponding map. Likewise,
\(J_n\) is dimensionless; identifying it with a physical Hamiltonian
requires the representation and the action and clock scale used by
that sector.
\end{falsifier}

\begin{theorem}[Faithfulness of the marked spectral realization]
\label{thm:jacobi-marked-faithfulness}
The infinite Jacobi operator \(J\), together with its unit cyclic
vector \(e_0\) and charge \(Q\), determines the register \(K\).
Two realizations of this type are unitarily equivalent through a map
preserving the marked vector and charge if and only if their registers
agree. The spectrum as a set, by contrast, is \([0,1]\) for every
positive register in the domain.
\end{theorem}
\begin{proof}
Polynomials are dense in \(L^2(\mu_K)\): first approximate continuous
functions and then use their density for the finite measure.
The constructed orthonormal basis defines a unitary map taking
\(e_n\) to \(\psi_n\). Under this map, \(J\) is multiplication
by \(t\), a bounded self-adjoint operator, and \(e_0\) represents
the constant function one. Therefore,
\[
 \langle e_0,J^re_0\rangle=m_r\qquad(r\geq0).
\]
These moments determine the measure: two measures with the same
moments agree on polynomials, then on continuous functions by uniform
approximation, and finally as Borel measures. Its distribution
\(F(t)=\mu_K([0,t])\) is continuous and strictly increasing.
Each initial interval has mass \(1/12\), so
\[
 t_j=F^{-1}(j/12),\qquad
 K_j=Q\bigl(F^{-1}(j/12)-F^{-1}((j-1)/12)\bigr).
\]
A unitary equivalence preserving \(e_0\) preserves the moments and
therefore the register. The converse follows from the unique
construction.

For the last statement, the density of \(\mu_K\) is positive on
every subinterval. Each \(\lambda\in[0,1]\) admits unit-norm
functions concentrated on progressively smaller intervals around
\(\lambda\), on which the norm of \((J-\lambda I)f\) tends
to zero. Thus the entire interval belongs to the spectrum. Outside
it, multiplication by \(1/(t-\lambda)\) is bounded and defines
the resolvent. This proves equality of the spectrum.
\end{proof}

The theorem distinguishes the extent of the spectrum from the
information it preserves relative to a marked observation. The same
spectral interval admits the different distributions of the collection;
the measure observed from the unit recovers the original separations.
The link with the Grassmannian representation is now reversible:
marked minors reconstruct \(K\), the register constructs
\((J,e_0,Q)\), and its moments and quantiles recover the same
register. Every reader factoring through \((K,\xi)\) therefore
has a marked spectral expression once the remaining coordinates
\(\xi\) are preserved. The spectral realization contains structure,
in addition to pole positions or eigenvalues.
